% !TeX root = ../main.tex
\section{Заключение}

В ходе выполнения лабораторной работы были экспериментально исследованы основные линейные и нелинейные параметры усилителя высокочастотного и СВЧ-диапазонов с применением современной контрольно-измерительной аппаратуры (цифрового анализатора спектра RIGOL DSA1030 и высокочастотных генераторов Advantex SG8).

\textbf{Основные результаты проведенных исследований:}
\begin{enumerate}
    \item \textbf{Линейные характеристики передачи ($S_{21}$) и обратной изоляции ($S_{12}$):} 
    Применение метода замещения (калибровки измерительного тракта без усилителя) позволило исключить частотную неравномерность генератора, измерительных кабелей и входных цепей анализатора. В диапазоне частот 10--2900~МГц исследуемый усилитель обеспечивает максимальный коэффициент усиления $S_{21} \approx 27\text{--}28\text{ дБ}$ на частотах 100--200~МГц с плавным спадом до $14.3\text{ дБ}$ на частоте 2900~МГц, обусловленным влиянием паразитных емкостей транзистора. Обратная развязка ($S_{12}$) ухудшается с увеличением частоты (от $-21.0\text{ дБ}$ до $-3.17\text{ дБ}$), что отражает рост проходной емкостной связи между выходом и входом каскада.

    \item \textbf{Шумовые параметры ($NF$):} 
    Косвенным методом измерения мощности теплового шума на выходе каскада при согласованной входной нагрузке 50~Ом получена частотная зависимость коэффициента шума с обязательным учетом полосы фильтра промежуточной частоты анализатора ($\text{RBW} = 100\text{ кГц}$, $50\text{ дБГц}$) и затухания в выходном соединительном кабеле. В диапазоне 500--1700~МГц коэффициент шума составляет $2.6\text{--}3.2\text{ дБ}$, монотонно возрастая до $4.12\text{ дБ}$ на частоте 2900~МГц. Полученные значения подтверждают применимость исследуемого усилителя в качестве входного малошумящего каскада (МШУ) радиоприемного тракта дециметрового диапазона волн.

    \item \textbf{Амплитудная характеристика и точка децибельной компрессии ($P_{1\text{dB}}$):} 
    На частоте $f = 1\text{ ГГц}$ экспериментально снята амплитудная характеристика каскада в диапазоне входных мощностей от $-10\text{ дБм}$ до $+18\text{ дБм}$. Определена точка компрессии усиления на 1~дБ: $IP_{1\text{dB}} = -1.36\text{ дБм}$, $OP_{1\text{dB}} = -0.64\text{ дБм}$. Данная точка устанавливает верхнюю границу линейного динамического диапазона усилителя, выше которой нарушается линейность передачи и начинается режим насыщения с предельной выходной мощностью $P_{\text{sat}} \approx +5\text{ дБм}$.

    \item \textbf{Нелинейность по интермодуляции ($IP_3$):} 
    Методом двухтонового воздействия ($f_1 = 999\text{ МГц}$, $f_2 = 1000\text{ МГц}$) исследованы комбинационные составляющие 3-го порядка ($2f_1 - f_2$ и $2f_2 - f_1$). В малосигнальной области подтверждено теоретическое соотношение наклонов характеристик: приращение входной мощности на 1~дБ вызывает рост мощности продуктов $IM_3$ на 3~дБ (наклон 3:1), что приводит к ухудшению параметра $IMD_3$ на 2~дБ. Графической и аналитической экстраполяцией характеристик найдены координаты точки пересечения третьего порядка: $IIP_3 = +20.0\text{ дБм}$, $OIP_3 = -1.0\text{ дБм}$. Превышение точки $IIP_3$ над точкой $IP_{1\text{dB}}$ более чем на 20~дБ подтверждает высокую динамическую линейность исследуемого каскада в штатном малосигнальном режиме приема.
\end{enumerate}

Таким образом, поставленные цели лабораторной работы достигнуты в полном объеме: освоены методы калибровки СВЧ-трактов, измерения волновых параметров рассеяния, шумовой меры и нелинейных искажений усилительных устройств.
